\documentclass{article} % This command is used to set the type of
% document you are working on such as an article, book, or presenation

\usepackage{geometry} % This package allows the editing of the page layout
\usepackage{amsmath}  % This package allows the use of a large range
% of mathematical formula, commands, and symbols
\usepackage{graphicx}  % This package allows the importing of images
\usepackage{
  amsmath,      % Math Environments
  amssymb,      % Extended Symbols
  enumerate,        % Enumerate Environments
  graphicx,      % Include Images
  lastpage,      % Reference Lastpage
  multicol,      % Use Multi-columns
  multirow      % Use Multi-rows
}

\newcommand{\question}[2][]{
  \begin{flushleft}
    \textbf{Question #1}: \textit{#2}

\end{flushleft}}
\newcommand{\sol}{\textbf{Solution}:} %Use if you want a boldface solution line
\newcommand{\maketitletwo}[2][]{
  \begin{center}
    \Large{\textbf{Assignment #1}

    MATH 4362} % Name of course here
    \vspace{5pt}

    \normalsize{N. Ohayon Rozanes  % Your name here

    \today}        % Change to due date if preferred
    \vspace{15pt}

\end{center}}
\begin{document}
\maketitletwo[5]  % Optional argument is assignment number
%Keep a blank space between maketitletwo and \question[1]

\question[1]{}
\begin{enumerate}
  \item $$f(x) = e^{-(x+4)^2}$$
    \begin{align*}
      \hat f(k) &= \frac{1}{\sqrt{2\pi}}\int_{-\infty}^\infty f(x)e^{-ikx} dx\\
      &= \frac{1}{\sqrt{2\pi}}\int_{-\infty}^\infty e^{-(x+4)^2}e^{-ikx} dx \\
      &= \frac{1}{\sqrt{2\pi}}\int_{-\infty}^\infty e^{-(x+4)^2 - ikx} dx
    \end{align*}

    Further analyzing the exponent:
    \begin{align*}
      -(x+4)^2 - ikx &= -x^2 - x(8+ ik) - 16 \\
      &=-(x+\frac{8+ik}{2})^2 + \frac{(8+ik)^2}{4} -16\\
    \end{align*}
    substituting back into the integral and pulling out terms:
    \begin{align*}
      \frac{1}{\sqrt{2\pi}}\int_{-\infty}^\infty e^{-(x+4)^2 - ikx}
      dx &= \frac{\exp(\frac{(8+ik)^2}{4} - 16)}{\sqrt{2\pi}}
      \int_{-\infty}^\infty \exp(-(x+\frac{8+ik}{2})^2) dx
    \end{align*}
    Which we recognize as the guassian integral, so the integral evaluates to
    \begin{align*}
      \hat f(k) = \frac{\exp(\frac{(8+ik)^2}{4} - 16)}{\sqrt{2}} =
      \frac{\exp(4ik -\frac{k^2}{4})}{\sqrt{2}}
    \end{align*}
  \item
    $$f(x) = e^{-|x+1|}$$
    Given that we know that for $g(x) = e^{-|x|}$ $$\hat g(x) =
    \sqrt{\frac{2}{\pi}}\frac{1}{k^2+1}$$,
    We observe that $g(x+1) = g(x-(-1)) = f(x)$, then applying the
    shift thereom,
    we have that:
    $$
    \hat f(x) = e^{ik}\hat g(k)  = \sqrt{\frac{2}{\pi}}\frac{e^{ik}}{k^2+1}
    $$
\end{enumerate}
\question[2]{}
\begin{enumerate}
  \item
    \[
      \hat f(k) = e^{-k^2}
    \]
    Then we have that
    \[
      f(x) = \frac{1}{\sqrt{2\pi}}\int_{-\infty}^\infty \hat f e^{ikx}dk
      = \frac{1}{\sqrt{2\pi}}\int_{-\infty}^\infty e^{-k^2 + ikx}dk
    \]
    Again, we can complete the square in the exponent:
    \begin{align*}
      -k^2 + ikx = -(k - \frac{ix}{2})^2 + \frac{x^2}{4}
    \end{align*}
    So the integral becomes recognizable as the Guassian integral
    \[
      \exp\left(\frac{x^2}{4}\right) \int_{-\infty}^\infty
      e^{-(k-\frac{ix}{2})^2} dk = \exp(\frac{x^2}{4})\sqrt{\pi}
    \]
    Then we have that
    $$
    f(x) = \frac{\exp\left(\frac{x^2}{4}\right)}{\sqrt{2}}
    $$
  \item
    \[
      \hat f(k) = e^{-|k|} = e^{-k}\sigma(k) + e^{k}\sigma(-k)
    \]
    As a linear operator, we consider the inverse fourier transform
    of each component seperatly:
    \begin{align*}
      \int_{-\infty}^\infty e^{-k}e^{ikx}\sigma(k) dk &=
      \int_0^\infty e^{-k(1 - ix)} \\
      &= \frac{1}{(1-ix)}e^{-k(1-ix)}\Big\vert_0^\infty \\
      &= \frac{1}{1-ix}
    \end{align*}
    Similarly
    \begin{align*}
      \int_{-\infty}^\infty e^{k}e^{ikx}\sigma(-k) dk &=
      \int_{-\infty}^0 e^{k(1+ix)} \\
      &= \frac{1}{1+ix}e^{k(i+ix)}\Big\vert_{-\infty}^0 \\
      &= \frac{1}{1+ix}
    \end{align*}
    Then we have that
    $$
    f(x) = \frac{1}{\sqrt{2\pi}}\left(\frac{1}{1+ix} + \frac{1}{1-ix}\right)
    $$
    Which we simplfy to
    \[
      f(x) = \frac{1}{\sqrt{2\pi}}\frac{2}{1+x^2} =
      \sqrt{\frac{2}{\pi}}\frac{1}{1+x^2}
    \]
\end{enumerate}
\question[3]{}
\begin{enumerate}
  \item
    \[
      \begin{cases}
        u_t = u_{xx} \\
        u(0,x) = e^{-x^2}
      \end{cases}
    \]
    We begin by finding the Fourier transform of the initial
    condition $f(x) = e^{-x^2}$, luckily, this is a known transform
    \[
      \hat f(k) = \frac{e^{-k^2/4}}{2\sqrt{2}}
    \]
    So we have the solution to the PDE in frequency:
    \[
      \hat u(t,k) = \frac{1}{2\sqrt{2}}e^{-k^2t}e^{-k^2/4} =
      \frac{1}{2\sqrt{2}}e^{-k^2(t + 1/4)}
    \]
    Taking the inverse fourier transform
    \begin{align*}
      u(t,k) &= \frac{1}{\sqrt{2\pi}}\int_{-\infty}^\infty
      \frac{1}{2\sqrt{2}}e^{-k^2(t + 1/4) + ikx} dk\\
      &= \frac{1}{4\sqrt{\pi}}\int_{-\infty}^\infty
      e^{-k^2(t + 1/4) + ikx}\\
    \end{align*}
    Completing the square on the exponent, it becomes:
    \[
      -\left(t+\frac{1}{4}\right)\left(k-\frac{ix}{2(t+\frac{1}{4})}\right)^2
      - \frac{x^2}{4(t + \frac{1}{4})}
    \]
    Thus the integral becomes
    \[
      e^{ - \frac{x^2}{4(t + \frac{1}{4})}}\int_{-\infty}^\infty
      \exp\left(-\left(t+\frac{1}{4}\right)\left(k-\frac{ix}{2(t+\frac{1}{4})}\right)^2\right)
    \]
    Which evaluates to
    \[
      e^{ - \frac{x^2}{4(t + \frac{1}{4})}}
      \sqrt{\frac{\pi}{(t+\frac{1}{4})}}
    \]
    Therefore, the solution to the PDE is:
    \[
      u(t,x) = \frac{e^{ - \frac{x^2}{4(t +
      \frac{1}{4})}}}{\sqrt{t+\frac{1}{4}}}
    \]
  \item For $f(x) = \sigma(x)$, we have the Fourier transform
    \[
      \hat f(k) = \sqrt{\frac{\pi}{2}}\delta(x) - \frac{i}{\sqrt{2\pi}k}
    \]
    then, in frequency, the solution is
    \[
      \hat u(t,k) = e^{-k^2t}\left(\sqrt{\frac{\pi}{2}}\delta(x) -
      \frac{i}{\sqrt{2\pi}k}\right)
    \]
    Taking the inverse Fourier transform:
    \begin{align*}
      \frac{1}{\sqrt{2\pi}}\int_{-\infty}^\infty
      e^{-k^2t}\left(\sqrt{\frac{\pi}{2}}\delta(x) -
      \frac{i}{\sqrt{2\pi}k}\right) dk &=
      \frac{1}{\sqrt{2\pi}}\int_{-\infty}^\infty
      e^{-k^2t}\sqrt{\frac{\pi}{2}}\delta(x) -
      e^{-k^2t}\frac{i}{\sqrt{2\pi}k} dk
    \end{align*}
    We evaluate each term individually: using the properties of the
    $\delta$ distribution, the first is simply
    \[
      \int_{-\infty}^\infty e^{-k^2t}\sqrt{\frac{\pi}{2}}\delta(x) =
      \sqrt{\frac{\pi}{2}}
    \]
    I don't believe the second term of the integrable has an
    elemntary solution, so the final answer, as far as I can evaluate is:
    \[
      u(t,x) = \frac{1}{\sqrt{2\pi}}\left(\sqrt{\frac{\pi}{2}} -
      \frac{i}{\sqrt{2\pi}}\int_{-\infty}^\infty \frac{e^{-k^2t}}{k}dk\right)
    \]

\end{enumerate}
\end{document}
